\documentclass[a4paper,12pt]{article}
\usepackage[utf8]{inputenc}
\usepackage{amsmath, amssymb}
\usepackage{geometry}
\usepackage{graphicx}
\geometry{margin=2.5cm}

\title{Esercitazione 1}
\author{Gabriele Perna}
\date{}

\begin{document}

\section{Esercizio 1}
Mostrare le componenti fortemente connesse del seguente grafo orientato.

\[
\begin{aligned}
V &= \{a,b,c,d,e,f,g,h,i,j,k,l,m,n,o,p,q\} \\
E &= \varnothing
\end{aligned}
\]

Guardando l'immagine, $n$ non può fare parte della componente perché è sorgente.
Non ci sono pozzi in questo grafo.

Ogni coppia di $x$ e $y$ esiste un walk da $x$ a $y$.  
Una componente non è componente se può essere espansa.

Per trovare le componenti bisogna fare ad occhio.

\section{Esercizio 2}
Sia $G = (V,E)$ un DAG, e sia $H = (V,E^{op})$ il grafo ottenuto invertendo le direzioni degli archi di $G$.
Per ognuno dei seguenti enunciati dire se è vero: in caso affermativo fornire una dimostrazione; in caso negativo fornire un controesempio.

\begin{enumerate}
    \item Ogni sorgente in $G$ è pozzo in $H$
    \item Ogni pozzo in $G$ è sorgente in $H$
    \item Esiste un walk da $x$ a $y$ in $G$ se e solo se esiste un walk da $x$ a $y$ in $H$
    \item Esiste un walk da $x$ a $y$ in $G$ se e solo se esiste un walk da $y$ a $x$ in $H$
    \item $H$ è un DAG
\end{enumerate}

\paragraph{Risposte:}

Per sostenere le nostre risposte prenderemo questo grafo:
\[
G = \big(\{A,B\}, \{(A,B)\}\big)
\]

\begin{enumerate}
    \item Vero, il ruolo delle sorgenti e dei pozzi si invertono.
    \item Vero, come sopra.
    \item Falso, gli archi si invertono e quindi\ldots
    \item Vero, gli archi si invertono.
    \item Vero, un ciclo se esiste si formerà di nuovo, ma al contrario. Analogamente, non si potranno formare cicli in un grafo opposto.
\end{enumerate}

\section{Esercizio 3}
Per ognuno dei seguenti enunciati dire se è vero: in caso affermativo
fornire una dimostrazione; in caso negativo fornire un controesempio.

\begin{enumerate}
    \item Per tutti i DAG $G = (V,E)$, $E \in Rel(V,V)$ è un ordinamento parziale.
    \item Per tutti i DAG $G = (V,E)$, $E \cup id_V \in Rel(V,V)$ è un ordinamento parziale.
    \item Per tutti i DAG $G = (V,E)$, $E^* \in Rel(V,V)$ è un ordinamento parziale.
\end{enumerate}

\paragraph{Risposte:}
\begin{enumerate}
    \item Falso, ogni insieme $E$ non contiene le coppie di nodi necessarie per essere riflessivo.
    Per esempio se $V = \{a,b\}$ ed $E = \{(a,b)\}$, allora $E$ non è riflessiva e quindi non è un ordinamento parziale.
    \item Falso, l'unione con l'identità di $V$ sicuramente risolve il problema della riflessività, ma non abbiamo ancora la transitività.
    Per esempio, se $V = \{a,b,c\}$ e $E \cup id_V = \{(a,a),(b,b),(c,c),(a,b),(b,c)\}$, allora non è transitiva.
    \item Vero, la stella di Kleene è una chiusura riflessiva e transitiva e un DAG non può mai avere coppie simmetriche o cicli e, dato che è stato chiuso riflessivamente, $E^*$ è un ordinamento parziale.
\end{enumerate}

\section{Esercizio 4}
Dimostrare il Teorema di Caratterizzazione per le relazioni anti-simmetriche:

Per tutti gli insiemi $A$, e tutte le relazioni $R \in Rel(A,A)$ vale che:
\[
R \text{ è anti-simmetrica se e solo se } R \cap R^{op} \subseteq id_A
\]

\paragraph{Risposta:}
Ricordiamo il significato di anti-simmetria:
una relazione è anti-simmetrica se per ogni coppia $(a,b)$ non esiste $(b,a)$ a meno che $b = a$.  
L'insieme risultante di $R \cap R^{op}$ conterrà solo coppie riflessive, con $a = b$, nel caso $R$ sia anti-simmetrica.

Nel caso $R$ non sia anti-simmetrica e quindi per almeno una coppia $(a,b)$ esiste una coppia $(b,a)$,in $R^{op}$ avremo:
\[
\begin{aligned}
(a,b),(b,a) &\in R\\
(b,a),(a,b) &\in R^{op}
\end{aligned}
\]

Questi, come possiamo vedere, sono elementi uguali e quindi compariranno nell'insieme intersezione, non convalidando la condizione di anti-simmetria.

\section{Esercizio 7}

Definiamo la funzione $mult2: \mathbb{N} \rightarrow \mathbb{N}$ induttivamente come:
\begin{itemize}
    \item Clausola Base: $mult2(0) = 0;$
    \item Clausola Induttiva: $mult2(n+1) = 2+mult2(n)$\\
\end{itemize}

Definiamo la sommatoria $\mathbb{N} \times Fun(\mathbb{N},\mathbb{N} \rightarrow \mathbb{N})$ induttivamente come:
\begin{itemize}
    \item Clausola Base: $sommatoria(0,f) = f(0) per ogni f \in fun(\mathbb{N},\mathbb{N})$
    \item Clausola Induttiva: $sommatoria(n+1,f) = f(n+1) + sommatoria(n,f) per ogni f \in Fun(\mathbb{N},\mathbb{N})$\\
\end{itemize}


Dimostrare per induzione i seguenti due enunciati:
\begin{enumerate}
    \item Per ogni $n,m, \in \mathbb{N}, mult2(n+m) = mult2(n) + mult2(m)$
    \item Per ogni $n \in \mathbb{N} e f \in Fun(\mathbb{N},\mathbb{N}), sommatoria (n,f;mult2) = mult2(sommatoria(n,f))$
\end{enumerate}
 
\textbf{Risposta:}

\end{document}
